\subsection{window-\texorpdfstring{$6$}{6} 的手征中心带、bulk 数域不可再分性与奇素响应线}\label{subsec:conclusion-window6-chiral-bulk-flag-oddprime-rigidity}

沿用小节 \ref{subsec:conclusion-window6-matrix-shadow-boundary-walsh-duality}、小节 \ref{subsec:conclusion-window6-center-grassmann-abelian-boundary-nonrecoverability}、小节 \ref{subsec:conclusion-window6-permutation-module-laplacian-reynolds} 与小节 \ref{subsec:conclusion-window6-geometry-spectrum-arithmetic-information-obstructions} 的记号。前述链条已经分别把 \(\sigma_{\mathrm{geo}}\) 的手征压缩、群胚矩阵包络的半单分层、stable-type 推前核的 \(1\oplus 20\) 有理谱结构以及中心 dyadic 模上的 Grassmann 塌缩固定为可复核的闭式。将这些接口并置之后，还可进一步抽取出如下几条终端后果。

\begin{theorem}[\texorpdfstring{\(\sigma_{\mathrm{geo}}\)}{sigma\_geo}-等变四块超立方体正规形]\label{thm:conclusion-window6-sigma-geo-fourblock-normal-form}
\leanverified{paper\_conclusion\_window6\_sigma\_geo\_fourblock\_normal\_form}
设
\[
\Omega_6=\{0,1\}^6\cong \{0,1\}^2\times \Omega_4,
\qquad
\Omega_4:=\{0,1\}^4,
\]
其中 \((a,b)\in\{0,1\}^2\) 对应第 \(1\) 与第 \(5\) 坐标，\(u\in\Omega_4\) 对应其余四个坐标。记 \(e_{ab,u}\) 为对应基向量，并定义
\[
z_u:=e_{00,u}-e_{11,u},
\qquad
c_u:=e_{01,u}-e_{10,u},
\]
\[
p_u:=e_{00,u}+e_{11,u},
\qquad
s_u:=e_{01,u}+e_{10,u},
\qquad
r_u^{\pm}:=p_u\pm s_u.
\]
再令
\[
W_\chi:=\langle z_u:\ u\in\Omega_4\rangle,
\qquad
W_{\mathrm{bal}}:=\langle c_u:\ u\in\Omega_4\rangle,
\]
\[
W_{\pm}:=\langle r_u^{\pm}:\ u\in\Omega_4\rangle.
\]
则有规范直和分解
\[
\CC^{\Omega_6}
=
W_\chi\oplus W_{\mathrm{bal}}\oplus W_+\oplus W_-,
\qquad
\dim W_\chi=\dim W_{\mathrm{bal}}=\dim W_+=\dim W_-=16,
\]
并且六维超立方体邻接算子 \(\mathsf A_{Q_6}\) 满足
\[
\mathsf A_{Q_6}\!\mid_{W_\chi}\cong \mathsf A_{Q_4},
\qquad
\mathsf A_{Q_6}\!\mid_{W_{\mathrm{bal}}}\cong \mathsf A_{Q_4},
\]
\[
\mathsf A_{Q_6}\!\mid_{W_+}\cong \mathsf A_{Q_4}+2I,
\qquad
\mathsf A_{Q_6}\!\mid_{W_-}\cong \mathsf A_{Q_4}-2I.
\]
此外，\(\sigma_{\mathrm{geo}}\) 在 \(W_\chi\) 上按 \((-1)\) 作用，而在 \(W_{\mathrm{bal}}\)、\(W_+\) 与 \(W_-\) 上按 \((+1)\) 作用。
\end{theorem}
\begin{proof}
由定理 \ref{thm:window6-chiral-compression-hypercube-adjacency} 的定义，\(\sigma_{\mathrm{geo}}=\sigma_{1,5}\) 在 \((a,b)\in\{0,1\}^2\) 上作用为
\[
00\longleftrightarrow 11,
\qquad
01\longmapsto 01,
\qquad
10\longmapsto 10.
\]
故 \(z_u\) 为反不变量，而 \(c_u,p_u,s_u\) 都为不变量；从而 \(W_\chi\subset (\CC^{\Omega_6})^{-}\)，其余三块都落在 \((\CC^{\Omega_6})^{+}\)。

再计算邻接作用。把 \(\mathsf A_{Q_6}\) 分解为“残余四坐标翻转”与“第 \(1,5\) 坐标翻转”两部分。若 \(u'\sim u\) 表示 \(Q_4\) 中相邻，则
\[
\mathsf A_{Q_6}e_{ab,u}
=
\sum_{u'\sim u} e_{ab,u'}
+\mathsf A_{Q_2}e_{ab,u},
\]
其中 \(\mathsf A_{Q_2}\) 是 \(\{00,01,10,11\}\) 上的二维超立方体邻接算子。直接计算得
\[
\mathsf A_{Q_2}(e_{00}-e_{11})=0,
\qquad
\mathsf A_{Q_2}(e_{01}-e_{10})=0,
\]
\[
\mathsf A_{Q_2}(e_{00}+e_{11})=2(e_{01}+e_{10}),
\qquad
\mathsf A_{Q_2}(e_{01}+e_{10})=2(e_{00}+e_{11}).
\]
因此
\[
\mathsf A_{Q_6}z_u=\sum_{u'\sim u}z_{u'},
\qquad
\mathsf A_{Q_6}c_u=\sum_{u'\sim u}c_{u'},
\]
\[
\mathsf A_{Q_6}p_u=\sum_{u'\sim u}p_{u'}+2s_u,
\qquad
\mathsf A_{Q_6}s_u=\sum_{u'\sim u}s_{u'}+2p_u.
\]
故
\[
\mathsf A_{Q_6}r_u^{\pm}
=
\sum_{u'\sim u}r_{u'}^{\pm}\pm 2r_u^{\pm},
\]
即在 \(W_\chi\)、\(W_{\mathrm{bal}}\)、\(W_\pm\) 上分别得到 \(\mathsf A_{Q_4}\)、\(\mathsf A_{Q_4}\)、\(\mathsf A_{Q_4}\pm 2I\)。

最后，四族向量分别由 \(\Omega_4\) 参数化，故各自维数都是 \(16\)；而它们正给出 \(\{0,1\}^2\) 上 \(Q_2\) 邻接算子的四个特征方向，因而在每个 \(u\) 纤维上构成一组基。于是总维数
\[
16+16+16+16=64=|\Omega_6|
\]
并且上述和为直和。证毕。
\end{proof}

\begin{corollary}[window-\texorpdfstring{$6$}{6} 的中心谱带双重化]\label{cor:conclusion-window6-centerband-doubling}
在定理 \ref{thm:conclusion-window6-sigma-geo-fourblock-normal-form} 的设定下，
\[
\chi_{Q_6}(\lambda)
=
\chi_{Q_4}(\lambda)^2\,
\chi_{Q_4}(\lambda-2)\,
\chi_{Q_4}(\lambda+2).
\]
特别地，\(\mathsf A_{Q_4}\) 对应的中心谱带在 \(Q_6\) 中精确以重数 \(2\) 出现；这两份拷贝分别由手征块 \(W_\chi\) 与平衡块 \(W_{\mathrm{bal}}\) 承载。
\end{corollary}
\begin{proof}
由定理 \ref{thm:conclusion-window6-sigma-geo-fourblock-normal-form}，
\[
\mathsf A_{Q_6}
\cong
\mathsf A_{Q_4}\oplus \mathsf A_{Q_4}\oplus (\mathsf A_{Q_4}+2I)\oplus (\mathsf A_{Q_4}-2I).
\]
对特征多项式取乘积即得所述公式。证毕。
\end{proof}

\begin{theorem}[stable-type 读出的手征失明]\label{thm:conclusion-window6-hidden-chiral-band-invisible}
设 \(F_6:\Omega_6\to X_6\) 为 stable-type 读出。若 \(g\in\CC^{\Omega_6}\) 仅依赖于 stable type，即存在 \(\bar g\in\CC^{X_6}\) 使
\[
g=\bar g\circ F_6,
\]
则对任意 \(v\in W_\chi\) 都有
\[
\langle g,v\rangle=0.
\]
换言之，定理 \ref{thm:conclusion-window6-sigma-geo-fourblock-normal-form} 中的手征中心带 \(W_\chi\) 在对象层可见，但对 stable-type 读出完全不可见。
\end{theorem}
\begin{proof}
由定理 \ref{thm:foldbin6-geo-stabilizer-z2}，\(\sigma_{\mathrm{geo}}\) 属于 \(F_6\) 的几何稳定子，故 \(F_6\circ \sigma_{\mathrm{geo}}=F_6\)。因此任意仅依赖 stable type 的 \(g\) 都满足
\[
\sigma_{\mathrm{geo}}\cdot g=g.
\]
另一方面，由定理 \ref{thm:conclusion-window6-sigma-geo-fourblock-normal-form}，\(W_\chi\) 正是 \(\sigma_{\mathrm{geo}}\) 的 \((-1)\)-等变块。于是对任意 \(v\in W_\chi\)，利用置换表示对计数内积的酉性，得
\[
\langle g,v\rangle
=
\langle \sigma_{\mathrm{geo}}\cdot g,\sigma_{\mathrm{geo}}\cdot v\rangle
=
\langle g,-v\rangle
=
-\langle g,v\rangle.
\]
故 \(\langle g,v\rangle=0\)。证毕。
\end{proof}

\begin{theorem}[groupoid involution 的矩阵 Lefschetz 恒等式]\label{thm:conclusion-window6-groupoid-involution-lefschetz}
令
\[
\mathcal A_6:=\CC[\mathcal R_6^{\mathrm{bin}}],
\]
并记 \(\alpha\) 为 \(\sigma_{\mathrm{geo}}\) 在 \(\mathcal A_6\) 上诱导的线性 involution。定义其 \(\pm1\)-本征空间
\[
\mathcal A_6^{\alpha,\pm}:=\{a\in \mathcal A_6:\ \alpha(a)=\pm a\}.
\]
则有
\[
\dim_{\CC}\mathcal A_6^{\alpha,+}=148,
\qquad
\dim_{\CC}\mathcal A_6^{\alpha,-}=64.
\]
特别地，
\[
\dim_{\CC}\mathcal A_6^{\alpha,-}=|\Omega_6|=2^6,
\qquad
\operatorname{tr}(\alpha\mid \mathcal A_6)=84=4|X_6|.
\]
\end{theorem}
\begin{proof}
由定理 \ref{thm:terminal-window6-geo-fixed-subalgebra-wedderburn}，
\[
\mathcal A_6^{\alpha,+}
=
\mathcal A_6^\alpha
\cong
\CC^{\oplus 14}\oplus M_2(\CC)^{\oplus 12}\oplus M_3(\CC)^{\oplus 6}\oplus M_4(\CC)^{\oplus 2},
\]
故
\[
\dim_{\CC}\mathcal A_6^{\alpha,+}
=
14+12\cdot 4+6\cdot 9+2\cdot 16
=
148.
\]
另一方面，由定理 \ref{thm:conclusion-window6-matrix-envelope-abelian-shadow}，
\[
\mathcal A_6
\cong
M_2(\CC)^{\oplus 8}\oplus M_3(\CC)^{\oplus 4}\oplus M_4(\CC)^{\oplus 9},
\]
从而
\[
\dim_{\CC}\mathcal A_6
=
8\cdot 4+4\cdot 9+9\cdot 16
=
212.
\]
故
\[
\dim_{\CC}\mathcal A_6^{\alpha,-}
=
212-148
=
64
=
|\Omega_6|.
\]
最后，
\[
\operatorname{tr}(\alpha\mid \mathcal A_6)
=
\dim \mathcal A_6^{\alpha,+}-\dim \mathcal A_6^{\alpha,-}
=
148-64
=
84.
\]
又 \(X_6\) 的基数为 \(21\)，故 \(84=4|X_6|\)。证毕。
\end{proof}

\begin{theorem}[stable-type 核的有理 étale 对易代数]\label{thm:conclusion-window6-stabletype-rational-etale-commutant}
令
\[
\mathcal C_{6,\QQ}
:=
\operatorname{Cent}_{M_{21}(\QQ)}(P_6).
\]
再记
\[
\chi_{P_6}(x)=(x-1)q_6(x),
\qquad
K_6:=\QQ[x]/(q_6).
\]
则有规范 \(\QQ\)-代数同构
\[
\mathcal C_{6,\QQ}
=
\QQ[P_6]
\cong
\QQ[x]\big/\bigl((x-1)q_6(x)\bigr)
\cong
\QQ\times K_6.
\]
其中 \(q_6\) 为次数 \(20\) 的不可约多项式，故 \(K_6/\QQ\) 为 \(20\) 次数域，并且不存在真中间域
\[
\QQ\subsetneq E\subsetneq K_6.
\]
因此，window-\(6\) 的 stable-type 非平凡谱并不分裂成若干较小有理块，而是压缩成单一的 \(20\) 次数域。
\end{theorem}
\begin{proof}
由定理 \ref{thm:terminal-windowm-rational-commutant-idempotent-rigidity}（取 \(m=6\)），已有
\[
\mathcal C_{6,\QQ}\cong \QQ\times K_6.
\]

另一方面，定理 \ref{thm:terminal-window6-pushforward-charpoly-galois} 给出 \(q_6\) 不可约且无重根，因此 \((x-1)q_6(x)\) 亦无重根。由于其次数为
\[
1+20=21=\dim_{\QQ}\QQ^{X_6},
\]
故 \(P_6\) 的最小多项式恰为 \((x-1)q_6(x)\)。于是
\[
\QQ[P_6]\cong \QQ[x]\big/\bigl((x-1)q_6(x)\bigr).
\]
该代数的 \(\QQ\)-维数为 \(21\)。而 \(\QQ\times K_6\) 的维数也是
\[
1+[K_6:\QQ]=1+20=21.
\]
由于 \(\QQ[P_6]\subseteq \mathcal C_{6,\QQ}\) 且两者维数相同，故
\[
\mathcal C_{6,\QQ}=\QQ[P_6].
\]
再由 \((x-1)\) 与 \(q_6(x)\) 互素，中国剩余定理给出
\[
\QQ[x]\big/\bigl((x-1)q_6(x)\bigr)\cong \QQ\times \QQ[x]/(q_6)=\QQ\times K_6.
\]
最后，“无真中间域”正是推论 \ref{cor:terminal-window6-q6-root-field-primitive-nonsolvable} 的内容。证毕。
\end{proof}

\begin{corollary}[非平凡 bulk 的有理不可再分性]\label{cor:conclusion-window6-rational-bulk-indecomposable}
沿用定理 \ref{thm:conclusion-window6-stabletype-rational-etale-commutant} 的记号，并记
\[
\Pi_{\mathrm{stat}}:=\Pi_6,
\qquad
\Pi_{\mathrm{bulk}}:=I-\Pi_6.
\]
则 \(\mathcal C_{6,\QQ}\) 的幂等元恰为
\[
0,\qquad I,\qquad \Pi_{\mathrm{stat}},\qquad \Pi_{\mathrm{bulk}}.
\]
特别地，不存在任何与 \(P_6\) 对易的精确有理投影把 \(20\) 维非平凡 bulk 再切成真非零直和。
\end{corollary}
\begin{proof}
由定理 \ref{thm:conclusion-window6-stabletype-rational-etale-commutant}，
\[
\mathcal C_{6,\QQ}\cong \QQ\times K_6.
\]
域中的幂等元只有 \(0\) 与 \(1\)，故直积环中的幂等元也只有
\[
(0,0),\quad (1,1),\quad (1,0),\quad (0,1).
\]
它们对应于 \(0\)、\(I\)、\(\Pi_{\mathrm{stat}}\) 与 \(\Pi_{\mathrm{bulk}}\)。因此除稳态线与其补空间外，再无精确有理投影。证毕。
\end{proof}

\begin{theorem}[window-\texorpdfstring{$6$}{6} 有理谱函数演算的四秩量子化]\label{thm:conclusion-window6-rational-functional-calculus-four-rank}
设
\[
V_{\QQ}:=\QQ\langle X_6\rangle,
\qquad
\chi_{P_6}(\lambda)=(\lambda-1)q_6(\lambda),
\qquad
\deg q_6=20,
\]
其中 \(q_6\) 在 \(\QQ\) 上不可约且无重根。则对任意 \(f\in\QQ[\lambda]\)，都有严格四值秩分类
\[
\operatorname{rank}f(P_6)\in\{0,1,20,21\}.
\]
更精确地，
\[
\operatorname{rank}f(P_6)=
\begin{cases}
21,& \gcd(f,(\lambda-1)q_6)=1,\\[2mm]
20,& (\lambda-1)\mid f,\ q_6\nmid f,\\[2mm]
1,& q_6\mid f,\ (\lambda-1)\nmid f,\\[2mm]
0,& (\lambda-1)q_6\mid f.
\end{cases}
\]
因此，\(\QQ[P_6]\) 并不是一个连续可调的有理滤波器族；它只能产生零算子、稳态线提取、稳态余维 \(1\) 的 bulk 提取，或处处可逆算子。
\end{theorem}
\begin{proof}
由定理 \ref{thm:conclusion-window6-stabletype-rational-etale-commutant}，
\[
\QQ[P_6]\cong \QQ[x]\big/\bigl((x-1)q_6(x)\bigr)\cong \QQ\times K_6,
\]
其中 \(K_6=\QQ[x]/(q_6)\) 为 \(20\) 次数域。
于是
\[
V_{\QQ}\cong \QQ[x]/(x-1)\ \oplus\ \QQ[x]/(q_6)
\]
在 \(\QQ[\lambda]\)-模意义下分裂为一维稳态块与二十维简单 bulk 块。

对任意 \(f\in\QQ[\lambda]\)，算子 \(f(P_6)\) 在第一块上按标量 \(f(1)\) 作用；在第二块上按 \(f\bmod q_6\) 作用。由于 \(q_6\) 不可约，\(\QQ[x]/(q_6)\) 是域，故 \(f\bmod q_6\) 要么为零、要么为可逆元。于是第二块的秩只能为 \(0\) 或 \(20\)，第一块的秩只能为 \(0\) 或 \(1\)。

因此 \(\operatorname{rank}f(P_6)\) 只能取四个可能值
\[
0,\ 1,\ 20,\ 21.
\]
并且：
\begin{enumerate}
  \item 若 \((\lambda-1)\nmid f\) 且 \(q_6\nmid f\)，则两块都可逆，故秩为 \(21\)；
  \item 若 \((\lambda-1)\mid f\) 但 \(q_6\nmid f\)，则稳态块消失而 bulk 块可逆，故秩为 \(20\)；
  \item 若 \(q_6\mid f\) 但 \((\lambda-1)\nmid f\)，则 bulk 块消失而稳态块非零，故秩为 \(1\)；
  \item 若 \((\lambda-1)q_6\mid f\)，则两块都消失，故秩为 \(0\)。
\end{enumerate}
这就得到题述分类。证毕。
\end{proof}

\begin{corollary}[window-\texorpdfstring{$6$}{6} 的 \texorpdfstring{$18\oplus3$}{18+3} 几何字典不是有理谱像]\label{cor:conclusion-window6-rational-functional-calculus-no-18plus3-image}
记
\[
V_{\mathrm{cyc}}:=\QQ\langle X_6^{\mathrm{cyc}}\rangle,
\qquad
V_{\mathrm{bdry}}:=\QQ\langle X_6^{\mathrm{bdry}}\rangle,
\]
其中
\[
\dim_{\QQ}V_{\mathrm{cyc}}=18,
\qquad
\dim_{\QQ}V_{\mathrm{bdry}}=3.
\]
则对任意 \(f\in\QQ[\lambda]\)，都不可能有
\[
\operatorname{im}f(P_6)=V_{\mathrm{cyc}}
\qquad\text{或}\qquad
\operatorname{im}f(P_6)=V_{\mathrm{bdry}}.
\]
更一般地，不存在任何来自 \(\QQ[P_6]\) 的有理谱分裂
\[
V_{\QQ}=W_{18}\oplus W_3,
\qquad
\dim_{\QQ}W_{18}=18,
\qquad
\dim_{\QQ}W_3=3.
\]
因此，window-\(6\) 的 \(18\oplus 3\) 根--boundary 字典只能作为外置几何分裂出现，而不能作为 \(P_6\) 的有理函数演算所切出的动力学频带。
\end{corollary}
\begin{proof}
由定理 \ref{thm:conclusion-window6-rational-functional-calculus-four-rank}，任意 \(f(P_6)\) 的像秩只能取
\[
0,\ 1,\ 20,\ 21.
\]
而
\[
\dim_{\QQ}V_{\mathrm{cyc}}=18,
\qquad
\dim_{\QQ}V_{\mathrm{bdry}}=3
\]
均不在该集合内，因此它们都不可能是任何 \(f(P_6)\) 的像空间。

若存在所述有理谱分裂，则投影到 \(W_{18}\) 与 \(W_3\) 的两个投影算子都属于 \(\QQ[P_6]\)，其像秩应分别为 \(18\) 与 \(3\)，这与上面的四秩分类矛盾。证毕。
\end{proof}

\begin{theorem}[手征中心带的 stable-type 不可回升性]\label{thm:conclusion-window6-chiral-band-nonrecoverability}
设 \(W_\chi\subset \CC^{\Omega_6}\) 为定理 \ref{thm:conclusion-window6-sigma-geo-fourblock-normal-form} 的手征块。则有如下严格不可逆性：
\begin{enumerate}
  \item 任何 stable-type 可观测量 \(g=\bar g\circ F_6\) 都与 \(W_\chi\) 正交；
  \item 任何仅由 \(P_6\) 的有理函数演算所产生的精确投影，都只能区分稳态线与全部非平凡 bulk，而不能在后者内部再精确切出真子块。
\end{enumerate}
因此，window-\(6\) 的 \(16\) 维手征中心带一旦向 stable-type 层塌缩，就在有理口径下不可回升；对象层中的几何手征扇区，在类型层最终只留下一个不可再分的 \(20\) 次数域。
\end{theorem}
\begin{proof}
第一条即定理 \ref{thm:conclusion-window6-hidden-chiral-band-invisible}。第二条即推论 \ref{cor:conclusion-window6-rational-bulk-indecomposable}。两者合并即得结论。证毕。
\end{proof}

\begin{theorem}[boundary 三维平面与 incident 旗标的外自同构消解]\label{thm:conclusion-window6-boundary-flag-outer-ambiguity}
记
\[
G_6^{\mathrm{bin}}:=\mathcal G_6^{\mathrm{bin}},
\qquad
A_{(2)}:=Z(G_6^{\mathrm{bin}})\cong (\ZZ/2\ZZ)^8.
\]
再由推论 \ref{cor:conclusion-window6-twofiber-boundary-internal-splitting} 记
\[
B_\partial:=E_\partial\le A_{(2)}
\]
为 boundary sheet-parity 对应的规范三维子空间。则 \(\Aut(G_6^{\mathrm{bin}})\) 在 \(A_{(2)}\) 上诱导整个 \(\GL(8,2)\) 的自然作用，因而在下列三个集合上都作用传递：
\[
\{\ell\le A_{(2)}:\ \dim \ell=1\},
\qquad
\{W\le A_{(2)}:\ \dim W=3\},
\]
\[
\{(\ell,W):\ \dim \ell=1,\ \dim W=3,\ \ell\subset W\}.
\]
特别地，\(B_\partial\) 不是 \(G_6^{\mathrm{bin}}\) 的 characteristic 子对象；并且对任意直线 \(\ell\le B_\partial\)，\(\ell\) 及 incident flag \((\ell\subset B_\partial)\) 也都不是纯群论的特征数据。
\end{theorem}
\begin{proof}
由定理 \ref{thm:terminal-window6-fold-gauge-aut-out-structure}，\(\Aut(G_6^{\mathrm{bin}})\) 在中心
\[
A_{(2)}=Z(G_6^{\mathrm{bin}})\cong \FF_2^8
\]
上诱导整个 \(\GL(8,2)\) 的自然线性作用。

\(\GL(8,2)\) 在非零向量集上的传递性给出对一维子空间的传递作用；对三维子空间的传递性则是 Grassmann 流形 \(\mathrm{Gr}(3,8;\FF_2)\) 的标准事实。对于 incident flag，取任意两组
\[
(\ell_1,W_1),\qquad (\ell_2,W_2),
\qquad
\dim \ell_i=1,\ \dim W_i=3,\ \ell_i\subset W_i.
\]
在 \(W_1,W_2\) 中分别取基
\[
W_1=\langle v_1,v_2,v_3\rangle,\qquad \ell_1=\langle v_1\rangle,
\]
\[
W_2=\langle w_1,w_2,w_3\rangle,\qquad \ell_2=\langle w_1\rangle,
\]
并把它们分别延拓为 \(\FF_2^8\) 的基。于是存在唯一的线性同构把 \(v_i\) 送到 \(w_i\)，故 \(\GL(8,2)\) 在 incident flag 集上亦作用传递。

因此，不可能有某个特定的三维子空间、直线或 incident flag 在全部自同构下保持不变。特别地，\(B_\partial\) 与任意 \(\ell\le B_\partial\) 及其旗标都不是 characteristic 数据。证毕。
\end{proof}

\begin{corollary}[dyadic incident 旗标的 \texorpdfstring{$680085$}{680085} 重代数等价性]\label{cor:conclusion-window6-680085-dyadic-flags}
在定理 \ref{thm:conclusion-window6-boundary-flag-outer-ambiguity} 的设定下，\(A_{(2)}\cong \FF_2^8\) 中满足
\[
\ell\subset W,
\qquad
\dim \ell=1,
\qquad
\dim W=3
\]
的 incident flag 总数恰为
\[
\frac{(2^8-1)(2^8-2)(2^8-2^2)}{(2^3-1)(2^3-2)(2^3-2^2)}
\cdot
\frac{2^3-1}{2-1}
=
97155\cdot 7
=
680085.
\]
并且这 \(680085\) 个 flag 在 \(\Aut(G_6^{\mathrm{bin}})\) 作用下属于同一轨道。
\end{corollary}
\begin{proof}
三维子空间的总数为 Gaussian 二项式
\[
\#\mathrm{Gr}(3,8;\FF_2)
=
\frac{(2^8-1)(2^8-2)(2^8-4)}{(2^3-1)(2^3-2)(2^3-4)}
=
97155.
\]
每个三维空间含有
\[
\frac{2^3-1}{2-1}=7
\]
条一维子空间，故 incident flag 总数为 \(97155\cdot 7=680085\)。
轨道唯一性则由定理 \ref{thm:conclusion-window6-boundary-flag-outer-ambiguity} 中的传递作用立即推出。证毕。
\end{proof}

\begin{theorem}[线性响应层的规范奇素线]\label{thm:conclusion-window6-canonical-oddprime-response-line}
令 \(M\in M_{21}(\ZZ)\) 为 window-\(6\) 的边耦合矩阵，并定义有限阿贝尔响应群
\[
\mathcal C_6:=\ZZ^{21}/M\ZZ^{21}.
\]
则有显式整因子分解
\[
\det(M)=-2^{12}\cdot 3^2\cdot 5\cdot 23.
\]
因此对任意素数
\[
p\notin\{2,3,5,23\}
\]
都有
\[
\ker(M\bmod p)=0.
\]
另一方面，在模 \(23\) 上存在规范的一维零响应线
\[
\ell_{23}^{(6)}:=\ker(M\bmod 23)\subset \FF_{23}^{21},
\qquad
\dim_{\FF_{23}}\ell_{23}^{(6)}=1.
\]
故 window-\(6\) 的响应层并非纯 dyadic；在奇素方向上，存在且仅存在一条由项目数据强制选出的规范 \(23\)-线。
\end{theorem}
\begin{proof}
由定理 \ref{thm:terminal-window6-fiber-edge-coupling-det}，
\[
\det(M)=-2^{12}\cdot 3^2\cdot 5\cdot 23.
\]
故当 \(p\notin\{2,3,5,23\}\) 时，\(\det(M)\not\equiv 0\pmod p\)，从而 \(M\bmod p\) 在 \(\FF_p\) 上可逆，即 \(\ker(M\bmod p)=0\)。这正是推论 \ref{cor:terminal-window6-fiber-edge-coupling-p-adic-window} 的内容。

对 \(p=23\)，定理 \ref{thm:terminal-window6-coupling-mod23-kernel-line} 给出
\[
\dim_{\FF_{23}}\ker(M\bmod 23)=1.
\]
故 \(\ell_{23}^{(6)}\) 为规范的一维子空间。证毕。
\end{proof}

\begin{conjecture}[window-\texorpdfstring{$6$}{6} 的最小刚化三联体]\label{conj:conclusion-window6-minimal-rigidifying-triad}
记
\[
\mathfrak F_6:=(\ell_\partial\subset B_\partial\subset A_{(2)}),
\qquad
\mathfrak K_6:=\QQ\times K_6,
\qquad
\mathfrak L_6:=\ell_{23}^{(6)},
\]
其中 \(B_\partial\) 为定理 \ref{thm:conclusion-window6-boundary-flag-outer-ambiguity} 的 boundary 三维平面，\(\ell_\partial\subset B_\partial\) 为由额外几何输入选定的一条直线，而 \(K_6\) 与 \(\ell_{23}^{(6)}\) 分别由定理 \ref{thm:conclusion-window6-stabletype-rational-etale-commutant} 与定理 \ref{thm:conclusion-window6-canonical-oddprime-response-line} 确定。
则猜想：window-\(6\) 的完整刚化并非由单一 Lie completion、单一谱域或单一信息量即可固定，而是由上述三联体最小地共同决定；更具体地，
\begin{enumerate}
  \item 删去 \(\mathfrak F_6\) 会重新出现 boundary 中心上的外自同构歧义；
  \item 删去 \(\mathfrak K_6\) 会重新出现 stable-type 非平凡 bulk 的有理不可分辨性；
  \item 删去 \(\mathfrak L_6\) 会重新出现奇素响应层的欠定性。
\end{enumerate}
因而，几何旗标、单一 \(20\) 次 bulk 数域与唯一模 \(23\) 响应线分别记录同一对象在 dyadic 中心、stable-type 谱与 odd-primary 响应层上的三种互不冗余刚化数据。
\end{conjecture}

\noindent
综上，window-\(6\) 的深层统一已不能由单一接口承载。微态超立方体在 \(\sigma_{\mathrm{geo}}\) 下精确裂成四个 \(Q_4\)-型谱块，其中手征中心带对 stable-type 读出完全失明；而 stable-type 层的有理交换子又把全部非平凡谱压缩成单一 \(20\) 次数域，只留下稳态/ bulk 二分。再往下，中心 dyadic 层虽然保留了规范的 boundary 三维平面，却无法从纯群论中继续选出任何内部旗标；与此同时，线性响应层在 dyadic 之外首先显露出的奇素刚性，恰以唯一的模 \(23\) 零响应线出现。故超立方体手征分裂、groupoid 矩阵 involution、stable-type bulk 数域、boundary 旗标与 odd-primary 响应线并非彼此并列的修辞，而是同一刚化问题在五个互不冗余方向上的严格投影。
\endinput
